\chapter{Общий родитель рождения ячеек, часов и энергии}
\label{ch:v10-cell-birth-clock-energy-common-parent-origin}

\section{Проверяемая гипотеза}

Предыдущая глава получила безразмерный прирост геометрической координаты
\begin{equation}
 \Delta\zeta=\frac13\log\frac{1+2x}{1+x},
 \qquad x=e^{-S_{\rm vac}}>0,
 \label{eq:v10-clock-delta-zeta}
\end{equation}
но оставила свободной длительность одного шага. Проверим, можно ли связать
рождение ячейки с автономными квантовыми часами одним положительным
родителем и тем самым получить физический темп без ручного задания связи.

\section{Минимальный резонансный носитель}

Возьмём двухуровневый счётчик рождения и двухуровневые часы. Используем
базис
\begin{equation}
 \{|0,0\rangle,|0,1\rangle,|1,0\rangle,|1,1\rangle\}.
\end{equation}
В нём безразмерный голый гамильтониан равен
\begin{equation}
 H_0=N_B\otimes I+I\otimes N_C
 =\operatorname{diag}(0,1,1,2),
 \qquad \operatorname{Spec}H_0=\{0^{(1)},1^{(2)},2^{(1)}\}.
 \label{eq:v10-clock-bare-hamiltonian}
\end{equation}
Две средние вершины образуют точную резонансную пару. Обмен
\begin{equation}
 G_B=|0,1\rangle\langle1,0|+|1,0\rangle\langle0,1|
 \label{eq:v10-clock-exchange}
\end{equation}
удовлетворяет
\begin{equation}
 [H_0,G_B]=0,
 \qquad \rank G_B=2,
 \qquad \operatorname{Spec}G_B=\{-1^{(1)},0^{(2)},1^{(1)}\}.
 \label{eq:v10-clock-exchange-data}
\end{equation}
Следовательно, рождение можно связать с часовым переходом без нарушения
голой суммы энергий. Это устанавливает допустимый носитель, но ещё не
выбирает размерный коэффициент гамильтониана.

\section{Общий родитель относительной связи}

Нормированная мера предыдущего гейта задаёт положительное число
\begin{equation}
 k_X:=\log\frac{1+2x}{1+x}=3\Delta\zeta>0.
 \label{eq:v10-clock-growth-coupling}
\end{equation}
Обозначим квадрат относительной силы обмена через $u=\chi^2$, а отношение
скорости рождения к частоте часов через
\begin{equation}
 \rho:=\frac{\Gamma_B}{\Omega},
 \qquad \Omega:=\frac{E_C}{\hbar}.
\end{equation}
Минимальный общий положительный родитель имеет вид
\begin{equation}
 \mathcal P_{BC}(u,\rho)
 =\frac12(u-k_X)^2+\frac12(\rho-u)^2.
 \label{eq:v10-clock-common-parent}
\end{equation}
Его единственный нуль и стационарная точка равны
\begin{equation}
 u=\rho=k_X.
 \label{eq:v10-clock-selected-point}
\end{equation}
Гессиан
\begin{equation}
 \operatorname{Hess}\mathcal P_{BC}
 =\begin{pmatrix}2&-1\\-1&1\end{pmatrix},
 \qquad \rank=2,
 \qquad \det=1,
 \label{eq:v10-clock-parent-hessian}
\end{equation}
имеет положительный спектр
$\{(3-\sqrt5)/2,(3+\sqrt5)/2\}$. Поэтому общий родитель действительно
совместно выбирает относительную силу обмена и относительный темп рождения.

\section{Полученный темп и оставшаяся орбита}

После размерной калибровки часами получаем
\begin{equation}
 \Gamma_B=k_X\frac{E_C}{\hbar},
 \qquad
 H_B:=\frac{\Gamma_B}{3}
 =\Delta\zeta\frac{E_C}{\hbar},
 \qquad
 \frac{H_B}{\Omega}=\Delta\zeta.
 \label{eq:v10-clock-relative-growth}
\end{equation}
Последнее отношение является слепым безразмерным следствием общей
конструкции. Однако преобразование
\begin{equation}
 E_C\longmapsto cE_C,
 \qquad t\longmapsto\frac{t}{c}
 \label{eq:v10-clock-scale-orbit}
\end{equation}
сохраняет фазу $\exp(-iE_Ct/\hbar)$ и все безразмерные формулы родителя.
Логарифмическая карта имеет одномерное ядро. Значит, общий родитель снял
относительную свободу, но не создал абсолютную энергию часов.

Чтобы дополнительно потребовать прежнюю условную амплитуду
$h_{\rm vac}=x/\sqrt{8\pi}$, пришлось бы отдельно задать
\begin{equation}
 \Omega_{\rm req}
 =\frac{3x}{\sqrt{8\pi}\log((1+2x)/(1+x))}.
 \label{eq:v10-clock-required-frequency}
\end{equation}
Эта частота не следует из построенного носителя или родителя и потому не
может считаться выведенной.

\begin{theorem}[Совместное происхождение относительного темпа]
Резонансный четырёхмерный носитель и родитель
$\mathcal P_{BC}$ совместно выбирают $\chi^2=\Gamma_B/\Omega=k_X$.
Отсюда строго следует безразмерное отношение
$H_B/\Omega=\Delta\zeta$.
\end{theorem}

\begin{nogocriterion}[Относительные часы не создают абсолютную энергию]
Если общий родитель зависит только от $\chi^2$, $\Gamma_B/\Omega$ и
нормированной меры рождения, то он инвариантен относительно совместного
масштабирования $(E_C,t)\mapsto(cE_C,t/c)$. Подстановка
$\Omega_{\rm req}$ без нового геометрического селектора является ручной
калибровкой, а не выводом космологического масштаба.
\end{nogocriterion}

\begin{protocol}[Статус общего родителя рождения и часов]\begin{enumerate}
\item резонансный энерго-сохраняющий носитель: \textbf{построен};
\item положительный общий родитель: \textbf{построен};
\item архитектура общего родителя: \textbf{$9/9$};
\item происхождение относительной связи: \textbf{$3/3$};
\item слепое отношение $H_B/\Omega$: \textbf{$\Delta\zeta$};
\item энергия часов и абсолютный темп: \textbf{$0/2$};
\item физическая космологическая постоянная: \textbf{$0/1$};
\item следующий гейт: \textbf{аудит геометрических кандидатов на
абсолютную энергию часов}.
\end{enumerate}\end{protocol}

Машинный сертификат сохранён в
\path{s2t_v10_cell_birth_clock_energy_common_parent_origin_gate_results.json}.

\begin{thebibliography}{9}
\bibitem{v10-clock-page-wootters}
D.~N.~Page, W.~K.~Wootters,
\emph{Evolution without evolution: Dynamics described by stationary
observables}, Physical Review D 27 (1983), 2885--2892.

\bibitem{v10-clock-erker}
P.~Erker, M.~T.~Mitchison, R.~Silva, M.~P.~Woods, N.~Brunner,
M.~Huber,
\emph{Autonomous quantum clocks: Does thermodynamics limit our ability to
measure time?}, Physical Review X 7 (2017), 031022.
\end{thebibliography}